% !PRINTED 394
\weberpage{394}
For let us put
\[
\begin{aligned}
N(\omega) &= a + 2^k b,\\
\omega_1 &= \omega (1 + 2^{k-1} x),
\end{aligned}
\]
where $x$ is rational. Then $2k - 2 > k$, and consequently
\[
\begin{aligned}
N(\omega_1) &\equiv (a + 2^k b)(1 + 2^k x) \ (\mathrm{mod}\ 2^{k+1});\\
&\equiv a + 2^k (b + a x);
\end{aligned}
\]
if therefore
\[
a x + b \equiv 0 \ (\mathrm{mod}\ 2),
\]
then
\[
N(\omega_1) \equiv a \ (\mathrm{mod}\ 2^{k+1}).
\]
From all this there now results the following theorem:

8. \gesperrt{If a number $r$ in $Z$ is a quadratic residue of every odd prime dividing $D$, and if moreover, for even $D$,}
\[
\begin{aligned}
& \text{a)} \left\{
\begin{array}{ll}
r \equiv 1 \ (\mathrm{mod}\ 8) & \text{if } \dfrac{D}{4} \equiv 0 \ (\mathrm{mod}\ 8),\\[4pt]
r \equiv +1, -1 \ (\mathrm{mod}\ 8) & \text{\glqq}\quad \dfrac{D}{4} \equiv 2 \ (\mathrm{mod}\ 8),\\[4pt]
r \equiv 1, 3 \ (\mathrm{mod}\ 8) & \text{\glqq}\quad \dfrac{D}{4} \equiv 6 \ (\mathrm{mod}\ 8),
\end{array}
\right. \\[6pt]
& \text{b)}\ r \equiv +1, -5 \ (\mathrm{mod}\ 4) \qquad \text{\glqq}\quad \dfrac{D}{4} \equiv 3, 4, 7 \ (\mathrm{mod}\ 8),
\end{aligned}
\]
\gesperrt{then $r$ is a norm residue of any arbitrary modulus $m$ coprime to $r$.}

These numbers $r$ form a group $R$, which we call the \gesperrt{group of absolute norm residues}.

If $b$ is a fixed quadratic non-residue of the odd prime $p$, and $z$ a number in $Z$, then either $z$ or $b^{-1} z$ is contained in the group of norm residues. So if $P$ is the group of norm residues of $p^k$, then
\[
\begin{aligned}
(Z, P) &= 1, \quad \text{when } p \text{ does not divide } D,\\
(Z, P) &= 2, \quad \text{when } p \text{ divides } D,
\end{aligned}
\]
further, if $L$ is the group of norm residues of $2^k$:
\[
\begin{aligned}
(Z, L) &= 1, & D &\equiv 1 \ (\mathrm{mod}\ 4),\\
& & D &\equiv 4, 20 \ (\mathrm{mod}\ 32),\\
(Z, L) &= 2, & D &\equiv 8, 12, 16, 24, 28 \ (\mathrm{mod}\ 32),\\
(Z, L) &= 4, & D &\equiv 0 \ (\mathrm{mod}\ 32)\,^{1}).
\end{aligned}
\]

\par\noindent\rule{4cm}{0.4pt}\par
$^{1}$) There is no need to fear that the symbol $(Z, P)$ for the index of a divisor of a group, explained in \S\,99 and in Bd.\,II, \S\,21, will be confused with the extended \gesperrt{Legendre-Jacobi} symbol similarly denoted in \S\,85.
